在上一讲中，我们讨论了静电平衡导体的性质：导体内部电场为零，导体本身是等势体，表面是等势面，
多余电荷只能分布在表面。对于一个已经达到静电平衡的导体，电势和净电荷是描述其状态的两个重要物理量。
若附近还有别的导体，一个导体所带电荷改变时，不仅自身电势会改变，还会通过静电感应改变其他导体表面的
电荷分布和电势。因此，导体组的静电问题本质上是在研究各导体的电荷、电势与几何结构之间的关系。

本讲先用电容描述这种关系，再讨论电容器的连接与储能，随后把结论推广到任意导体组，最后研究电介质在
电场中的极化。以下定义都默认导体的形状、相对位置及周围介质已经固定；只有在这些条件不变时，电荷与
电势之间的线性关系才可以由一个固定的电容或一组电容系数表示。

\subsection*{孤立导体的电容}

设空间中有若干导体，除其中一个以外，其余导体都接地，电势保持为零。使唯一未接地的导体带上净电荷
$Q$，静电平衡后它的电势为 $\varphi$。由静电场的线性叠加和唯一性定理可知：如果把 $Q$ 同时放大
$k$ 倍，空间中全部电场和电势也放大 $k$ 倍，因而 $Q/\varphi$ 不随电荷量改变。定义该导体在当前
几何环境中的单极电容为
\begin{equation}
  \boxed{C=\frac{Q}{\varphi}},\qquad Q=C\varphi.
\end{equation}
这里的“孤立”不是说宇宙中只有这一个导体，而是说只把它作为非零电势的导体，其余导体接地。电容取决于
导体的大小、形状、位置及周围介质，却与它在给定线性范围内带了多少电荷无关。国际单位制中电容的单位是
法拉（F），名称来自 Faraday；实际导体的电容通常远小于一法拉，常用 $\si{\micro F}$、$\si{nF}$、
$\si{pF}$。

半径为 $R$ 的孤立导体球是最容易计算的例子。球外电势与位于球心的点电荷相同，取无穷远电势为零，
\begin{equation}
  \varphi(R)=\frac{Q}{4\pi\varepsilon_0R},
  \qquad
  \boxed{C_{\rm sphere}=4\pi\varepsilon_0R}.
\end{equation}
结果与球所带电量无关，并且正比于球半径。球越大，在升高到同一电势时能够容纳的电荷越多。

给一个原来不带电的孤立导体逐渐充电。当它已经带有电荷 $q$ 时，再搬入微小电荷 $\dd q$ 所做的功是
$\varphi(q)\dd q=(q/C)\dd q$。从零充到 $Q$，系统的静电能为
\begin{equation}
  W=\int_0^Q\frac{q}{C}\dd q
  =\boxed{\frac{Q^2}{2C}=\frac12Q\varphi=\frac12C\varphi^2}.
\end{equation}
这些能量并不是“藏在电荷里”，而是以电场能的形式分布在空间中。真空静电场的能量密度为
$u_E=\varepsilon_0E^2/2\geq0$，所以静电系统的场能不会为负。

\subsection*{双极电容器与典型几何}

实际储存电荷时更常用两块彼此绝缘的导体。令两导体分别带 $+Q$ 与 $-Q$，电势为
$\varphi_+$、$\varphi_-$，两极电势差
$U=\varphi_+-\varphi_-$。只要几何与介质不变，$Q$ 与 $U$ 成正比，双极电容器的电容定义为
\begin{equation}
  \boxed{C=\frac{Q}{U}},\qquad Q=CU.
\end{equation}
孤立导体的 $Q/\varphi$ 与双极电容器的 $Q/U$ 形式相同，但物理边界条件不同，使用时不能混淆。

\arbitraryCapacitorDiagram

电容器不要求极板平行，也不要求形状规则。它的电容由两极的大小、形状、相对位置及介质共同决定。
当导体形状具有足够高的对称性时，可以先用 Gauss 定律求电场，再积分求电势差，最后取 $C=Q/U$。

对于面积为 $S$、间距为 $d$ 的平行板电容器，若极板线度远大于 $d$，中部电场近似均匀：
\begin{equation}
  E=\frac{Q}{\varepsilon_0S},\qquad
  U=Ed=\frac{Qd}{\varepsilon_0S},
  \qquad
  \boxed{C=\frac{\varepsilon_0S}{d}}.
\end{equation}
真实极板边缘的电场会弯曲，称为边缘效应。极板越大、间距越小，均匀场区域在全部场区中所占比例越大，
理想公式越准确。边缘场使实际电容通常略大于忽略边缘效应得到的数值。

同心球形电容器的内、外半径为 $R_1<R_2$。在两球面之间
\begin{equation}
  E(r)=\frac{Q}{4\pi\varepsilon_0r^2},
  \qquad
  U=\int_{R_1}^{R_2}E(r)\dd r,
\end{equation}
因此
\begin{equation}
  \boxed{C_{\rm sph}=4\pi\varepsilon_0\frac{R_1R_2}{R_2-R_1}}.
\end{equation}
令 $R_2\to\infty$，便退化为半径 $R_1$ 的孤立导体球。

同轴圆柱电容器内、外半径为 $a<b$，有效长度为 $l$，并且 $l$ 足够大，可忽略两端效应。若内柱带
$+Q$，两极之间距轴线 $r$ 处的场强为
\begin{equation}
  E(r)=\frac{Q}{2\pi\varepsilon_0lr}.
\end{equation}
由
\begin{equation}
  U=\int_a^bE(r)\dd r
   =\frac{Q}{2\pi\varepsilon_0l}\ln\frac ba
\end{equation}
得到
\begin{equation}
  \boxed{C_{\rm cyl}=\frac{2\pi\varepsilon_0l}{\ln(b/a)}}.
\end{equation}
圆柱电容正比于长度。外半径固定而内半径减小时，靠近细内导体处的场会迅速增强，这也说明“无限细带电线”
只是理论极限，真实导线半径会影响电势差和电容。

\capacitorGeometryDiagram

\subsection*{电容器的连接}

多个电容器接入电路后，可以用一个等效电容描述从外部端口看到的储荷效果。并联电容器的两端分别连接到
同一对节点，故各电容器的电势差都等于 $U$。第 $i$ 个电容器带电 $Q_i=C_iU$，电源搬运的总电荷为
\begin{equation}
  Q=\sum_iQ_i=\left(\sum_iC_i\right)U,
\end{equation}
所以
\begin{equation}
  \boxed{C_{\parallel}=\sum_iC_i}.
\end{equation}

串联电容器的中间连接导体与外界绝缘。若它初始不带净电荷，充电过程中电荷只能在这个导体内部重新分布，
净电荷始终为零，因而各串联电容器所带电荷量的绝对值相同。总电势差为
\begin{equation}
  U=\sum_iU_i=Q\sum_i\frac1{C_i},
\end{equation}
于是
\begin{equation}
  \boxed{\frac1{C_{\rm series}}=\sum_i\frac1{C_i}}.
\end{equation}
并联使等效电容增大，串联使等效电容减小。电容的串并联规律与电阻恰好相反，却与理想弹簧的等效关系在
形式上相似；真正的判断依据是各元件“电压相同”还是“电荷相同”。

\seriesParallelDiagram

\subsection*{电容储能与能量密度}

容量为 $C$ 的电容器最初不带电。当两极已经带有 $\pm q$ 时，电势差为 $u=q/C$，再把 $\dd q$ 从
负极搬到正极需要外力做功
\begin{equation}
  \dd W=u\dd q=\frac qC\dd q.
\end{equation}
从零充到 $Q$ 后，电容器储存的电场能为
\begin{equation}
  \boxed{W=\frac{Q^2}{2C}=\frac12QU=\frac12CU^2}.
\end{equation}
这个推导不依赖极板形状。对于平行板电容器，把上式除以板间体积 $Sd$，可得真空中的能量密度
\begin{equation}
  \boxed{u_E=\frac12\varepsilon_0E^2}.
\end{equation}
在线性、各向同性介质中，能量密度推广为 $u_E=\vect E\cdot\vect D/2=\varepsilon E^2/2$。

\subsection*{任意导体组与 Green 倒易定理}

孤立导体和双极电容器都只是导体组的特殊情况。设同一组固定导体有两个静电平衡状态：第一状态中第 $i$
个导体的电荷与电势为 $Q_i,\varphi_i$，第二状态为 $Q_i',\varphi_i'$。由 Green 第二恒等式可得
\begin{equation}
  \boxed{\sum_iQ_i\varphi_i'=\sum_iQ_i'\varphi_i}.
\end{equation}
它说明：把第一状态各导体的电荷乘以第二状态相应导体的电势后求和，与交换两个状态所得结果相同。
具有这种“交换激励与响应而结果不变”结构的结论称为倒易定理。它与静电场方程的线性、边界条件的唯一性
以及场能的二次型结构密切相关。

\reciprocitySphereDiagram

例如，取一个导体球和球外距球心为 $r$ 的点 $P$。状态 I 中球带单位电荷，$P$ 点无点电荷；球面电荷
均匀分布，它在 $P$ 点产生电势 $1/(4\pi\varepsilon_0r)$。状态 II 中球原来不带电，在 $P$ 点放置
电荷 $q$，静电感应使球面电荷重新分布，球仍为等势体，设其电势为 $\varphi$。对这两个状态应用倒易定理，
可立即得到
\begin{equation}
  \boxed{\varphi=\frac{q}{4\pi\varepsilon_0r}}.
\end{equation}
虽然状态 II 的球面电荷分布并不均匀，但球的电势无需逐点积分就能确定。

\reciprocityConductorsDiagram

再考虑两个任意形状导体。两种状态分别记为
$(Q_1,Q_2;\varphi_1,\varphi_2)$ 与
$(Q_1',Q_2';\varphi_1',\varphi_2')$，则
\begin{equation}
  Q_1\varphi_1'+Q_2\varphi_2'
  =Q_1'\varphi_1+Q_2'\varphi_2.
\end{equation}
由此可以导出多种有用结论。例如，在两种实验中轮流令一个导体带相同电荷、另一个保持中性，则一个导体
在另一个导体处产生的电势具有交换对称性；轮流令一个导体接地、另一个处于相同电势，接地导体上的总感应
电荷也具有交换对称性。把电荷写成各导体电势的线性函数
\begin{equation}
  Q_i=\sum_jC_{ij}\varphi_j,
\end{equation}
倒易定理给出
\begin{equation}
  \boxed{C_{ij}=C_{ji}}.
\end{equation}
因此，多导体静电感应并不是每一对导体各有一套互不相关的经验参数，而可以用一个对称电容系数矩阵统一描述。

\subsection*{电介质与相对介电常数}

真空电容器的两极之间没有物质。实际电容器往往在电极间填入橡胶、玻璃、纸、石英等绝缘材料，这类物质
称为电介质。实验表明：用某种均匀电介质完全填满电容器后，电容总会增大；更重要的是，对同一种介质，
无论原电容器形状如何，在适用的线性范围内，电容都近似按同一倍数增大。这个倍数称为相对介电常数
$\varepsilon_r$：
\begin{equation}
  \boxed{C=\varepsilon_rC_0},
\end{equation}
其中 $C_0$ 是相同几何结构在真空中的电容。常见被动介质的 $\varepsilon_r>1$。

\dielectricExperimentDiagram

\begin{center}
\small
\begin{tabular}{lcc lcc}
\toprule
材料 & 相对介电常数 $\varepsilon_r$ & 击穿场强/$(\si{V.m^{-1}})$ &
材料 & 相对介电常数 $\varepsilon_r$ & 击穿场强/$(\si{V.m^{-1}})$\\
\midrule
干燥空气 & $1.00059$ & $3\times10^6$ & 水 & $80$ & --\\
石英 & $3.78$ & $24\times10^6$ & 橡胶 & $6.7$ & $12\times10^6$\\
纸 & $3.7$ & $16\times10^6$ & 尼龙 & $3.4$ & $14\times10^6$\\
玻璃 & $5.6$ & $14\times10^6$ & 聚苯乙烯 & $2.56$ & $24\times10^6$\\
\bottomrule
\end{tabular}

{\small\color{TextGray}表 32.1\quad 室温下典型介质的相对介电常数和击穿电场（典型值）}
\end{center}

表中的击穿场强表示介质仍能保持绝缘的典型上限。超过这一数量级后，束缚电荷可能被大量释放，介质发生
电击穿，不能再用线性绝缘体模型处理。材料的介电常数与击穿场强还会受到温度、频率、纯度和制备状态影响，
表中数值只适合作数量级估计。

若保持电容器自由电荷 $Q$ 不变，完全插入介质后 $C$ 增大为 $\varepsilon_rC_0$，电势差
$U=Q/C$ 降低到原来的 $1/\varepsilon_r$，场强也相应减小。这说明介质内部出现了一个通常与外场反向的
附加电场。它来源于介质中束缚正、负电荷在外场下的微小相对位移。

\subsection*{介质极化与电位移}

经典图像中，绝缘体由大量整体中性的原子或分子组成，电子不能像金属自由电子那样在整个物体内迁移。
外电场作用时，原子中正、负电荷中心发生微小分离，形成诱导电偶极子；本来具有永久偶极矩的分子也会在
外场下出现取向偏好。两种机制都使大量微观偶极矩在统计平均后产生宏观极化。

单位体积内的总电偶极矩定义为极化强度
\begin{equation}
  \boxed{\vect P=\frac{\dd\vect p}{\dd V}}.
\end{equation}
若每个分子的平均偶极矩为 $\vect p$、数密度为 $N$，均匀极化时 $\vect P=N\vect p$。介质内部相邻偶极子
的正负电荷大体抵消，但在边界上会留下束缚面电荷；非均匀极化还会产生束缚体电荷：
\begin{equation}
  \boxed{\sigma_b=\vect P\cdot\uvec n},
  \qquad
  \boxed{\rho_b=-\nabla\cdot\vect P}.
\end{equation}

\polarizationDiagram

自由电荷与束缚电荷都产生真实电场。为了在介质问题中把由电极等带来的自由电荷单独列出，定义电位移矢量
\begin{equation}
  \boxed{\vect D=\varepsilon_0\vect E+\vect P}.
\end{equation}
于是 Gauss 定律可写为
\begin{equation}
  \boxed{\oint_S\vect D\cdot\dd\vect S=Q_{\rm free}},
  \qquad
  \nabla\cdot\vect D=\rho_{\rm free}.
\end{equation}
右侧只计自由电荷，极化形成的束缚电荷已经包含在 $\vect P$ 中。$\vect D$ 是重新组织电荷贡献后引入的
辅助场量，并不是另一种独立于电场的力。

“电位移”这个名称保留了电磁学发展早期的历史痕迹。十九世纪的机械以太模型曾把电场类比为某种连续介质
受力后的形变，把真空响应与普通物质的极化相加。以太图像后来被放弃，但由线性场方程整理出的
$\vect D=\varepsilon_0\vect E+\vect P$ 仍然是处理介质问题的有效定义。学习这一历史背景的意义在于区分：
公式可以保留，而最初解释公式的机械模型未必正确。

当场强不太大、介质各向同性且响应近似即时，极化强度可近似与电场成正比：
\begin{equation}
  \vect P=\varepsilon_0\chi_e\vect E.
\end{equation}
因此
\begin{equation}
  \boxed{\vect D=\varepsilon\vect E
  =\varepsilon_r\varepsilon_0\vect E},
  \qquad \varepsilon_r=1+\chi_e.
\end{equation}
从微观角度，若单个原子的诱导偶极矩满足 $\vect p=\alpha\vect E_{\rm loc}$，$\alpha$ 称为原子极化率。
宏观介电常数与 $N\alpha$ 有关；在最简单的各向同性模型中还要考虑邻近偶极子对局域场的修正，因此不能
总把 $\vect E_{\rm loc}$ 直接等同于宏观平均场 $\vect E$。测量介电常数曾是由宏观实验反推原子结构的重要途径。

\subsection*{介质边界、部分填充与受力}

在两种介质的交界面上作极薄柱形 Gauss 面和极窄矩形回路，可得静电边界条件
\begin{equation}
  \boxed{D_{2n}-D_{1n}=\sigma_f},
  \qquad
  \boxed{E_{2t}=E_{1t}}.
\end{equation}
只有自由面电荷使 $\vect D$ 的法向分量发生跳变；静电场的切向分量连续。界面没有自由电荷时
$D_{1n}=D_{2n}$，但若两侧介电常数不同，$E_n$ 一般不连续，电场线通过界面时会改变方向。

介质沿电场方向分层时，各层自由电荷面密度相同，$D_n$ 相同，而总电压等于各层电压之和，等效为电容串联：
\begin{equation}
  \frac1C=\frac{d_1}{\varepsilon_1S}+\frac{d_2}{\varepsilon_2S}+\cdots.
\end{equation}
介质若横向分割极板面积，各部分电势差相同、总电荷相加，等效为并联：
\begin{equation}
  C=\frac{\varepsilon_1S_1}{d}+\frac{\varepsilon_2S_2}{d}+\cdots.
\end{equation}
判断串联或并联不能只看图形，而要看各部分是电荷相同还是电压相同。

介质板插入电容器通常受到向内部的力，因为插入会使电容增大。电容器孤立、自由电荷固定时
\begin{equation}
  W_Q=\frac{Q^2}{2C},
  \qquad
  F_x=-\left(\frac{\partial W_Q}{\partial x}\right)_Q
  =\frac{Q^2}{2C^2}\frac{\partial C}{\partial x}.
\end{equation}
若电容器接恒压电源，必须把电源做功一起计入。电容器与电源合成系统的有效势能变化为
$\dd(W_C-W_s)=-U^2\dd C/2$，所以
\begin{equation}
  \boxed{F_x=\frac12U^2\frac{\partial C}{\partial x}}.
\end{equation}
两种边界条件给出的力方向一致，但推导中保持不变的量不同。

\subsection*{充电过程中的能量去向}

把未充电电容 $C$ 通过电阻 $R$ 接到理想电压源 $U_0$，最终电源搬运电荷 $Q=CU_0$，做功
$W_s=U_0Q=CU_0^2$。电容最终储能只有 $CU_0^2/2$，另一半在电阻中转化为热：
\begin{equation}
  Q_R=\int_0^\infty I^2R\dd t=\frac12CU_0^2.
\end{equation}
这一半能量与 $R$ 的具体数值无关，$R$ 只改变充电快慢。理想化地令 $R\to0$ 时，集总模型出现无限大
初始电流；真实电路中的寄生电感、电磁辐射和有限传播速度会参与能量转移，不能据此断言能量损耗消失。
串联电阻充电时
\begin{equation}
  U_C(t)=U_0\left(1-\e^{-t/(RC)}\right),
  \qquad
  I(t)=\frac{U_0}{R}\e^{-t/(RC)}.
\end{equation}
$RC$ 是时间常数。经过一个时间常数，电容电压达到终值的 $1-\e^{-1}$，电流降为初值的 $\e^{-1}$。

\subsection*{从边值问题理解电容}

电容的定义看上去只是一个比值，但它实际上概括了完整的静电边值问题。以孤立导体为例，给定导体表面的
电势以后，导体外部的电势必须满足 Laplace 方程，并在导体表面取给定常数、在无穷远趋于零。唯一性定理
保证满足这些条件的解只有一个。若把表面电势由 $\varphi$ 改成 $k\varphi$，原解整体乘以 $k$ 就满足新的
边界条件；于是电场、表面电荷密度和总电荷也全部乘以 $k$。这正是 $Q=C\varphi$ 成立的根本原因。

在这个构造中，其他接地导体虽然电势为零，却不等于“没有作用”。它们表面会出现感应电荷并改变空间电场，
因而也会改变未接地导体的电容。若未接地导体带正电，邻近接地导体朝向它的一侧通常出现负感应电荷。
接地意味着大地可以供给或吸收电子，所以这些导体的净电荷不必保持为零。把接地导体移近，通常会使给定
电荷造成的电势降低，从而使 $C=Q/\varphi$ 增大。这说明“一个导体的电容”仍然包含周围边界的信息。

双极电容器也由同样的线性原理决定。两极带等量异号电荷时，正极发出的电场线最终终止在负极，理想情况下
电场主要局限在两极之间。电容器储存的不是“正电荷”或“负电荷”本身，而是把正负电荷分离后形成的电场能。
如果保持两极形状不变，只把它们靠近，给定 $Q$ 造成的电势差会减小，因此电容增大；如果增大正对面积，
给定总电荷分摊到更大表面上，场强与电势差降低，电容同样增大。平行板公式中的 $S$ 和 $d$ 依赖关系由此
获得直观解释，而不只是代数结果。

平行板近似还包含一个容易被忽略的尺度条件。若极板边长与间距同量级，边缘弯曲的电场占据很大比例，不能
把全部场区都视为匀强场；此时 $C=\varepsilon_0S/d$ 只能给出粗略数量级。将极板边长同时放大而保持间距
不变，边缘区域的宽度仍约为 $d$，但中部匀强区域所占比例越来越大，近似才逐渐变好。竞赛题若强调
$l\gg d$，就是在要求忽略这部分端部修正。

圆柱电容的积分也展示了几何如何进入电容。圆柱 Gauss 面的面积正比于 $r$，所以场强按 $1/r$ 衰减；
电势差是 $1/r$ 从 $a$ 到 $b$ 的积分，因而出现 $\ln(b/a)$。如果只记结论而忽略这一过程，容易把它与球形
电容的 $1/R_1-1/R_2$ 混淆。三种典型电容器可用统一步骤处理：确定对称性，选取 Gauss 面求 $E$，沿场线
积分求 $U$，最后用 $C=Q/U$ 消去 $Q$。

\subsection*{连接、充电和储能的完整推理}

并联电容的关键是端点相同。达到静电平衡后，理想导线上的电势处处相同，因此每只电容的两端电压都等于
外部端口电压。电源向整个并联系统搬运的电荷是各电容正极板电荷之和，故等效电容相加。串联电容的关键
则是中间节点孤立：这个节点由相邻两块极板和连接导线构成，外部没有导电通路。它可以在内部重新分配正负
电荷，却不能改变净电荷；初始中性时，相邻两电容上的电荷量绝对值必然相同。把每段电压 $Q/C_i$ 相加，
便得到倒数相加的规律。

这种推理比直接背公式更可靠。例如串联电容的中间节点若预先带有净电荷，各电容所带电荷量便不再简单相等；
若电容器之间不是理想导线而连接了别的元件，节点条件也要重新分析。所谓“串联”或“并联”并不是由图形
看起来像一条线还是两条支路决定，而是由电势约束和电荷守恒决定。

电容储能中的因子 $1/2$ 来源于充电过程中电压从零逐渐增加。若错误地把最终电压 $U$ 乘以全部电荷 $Q$，
会得到 $QU$，比真实储能大一倍。正确做法是把每一小份电荷在当时电压 $u=q/C$ 下搬运所需的功积分起来。
也可以在 $q-u$ 图上理解：充电过程对应从原点到 $(Q,U)$ 的直线，储能等于直线下的三角形面积
$QU/2$。对于线性介质，$Q$ 与 $U$ 成正比；若介质具有强非线性或滞后，储能必须沿真实充电路径积分，
不能直接套用 $CU^2/2$。

把电场能密度在空间积分，应当得到与充电做功相同的总能量。以平行板电容为例，板间场近似均匀，
\[
 \int u_E\dd V=\frac12\varepsilon_0E^2Sd
 =\frac12\frac{\varepsilon_0S}{d}(Ed)^2
 =\frac12CU^2.
\]
这一等价性说明，集中参数 $C$ 的储能公式与分布在空间中的场能描述的是同一个物理事实。遇到介质只填充
部分空间、场分布不均匀的情形，直接积分场能有时比先求等效电容更方便。

电源给电容充电时还要区分“电源做功”和“电容最终储能”。理想恒压源每搬运 $\dd q$ 都做功
$U_0\dd q$，所以总功为 $U_0Q$；而电容电压从零升到 $U_0$，平均电压只有 $U_0/2$，最终储能是
$U_0Q/2$。差额必须进入电阻发热、辐射或电磁场的暂态过程。把电阻取得很小只会让过程更快，不会凭空改变
初态与终态之间的能量差。

\subsection*{倒易定理的三种典型用法}

Green 倒易定理的力量在于，它允许把一个难求的响应换成另一个容易计算的响应。应用时必须保证两个状态中
导体的形状、位置和介质完全相同，仅改变各导体的电荷或电势。若几何也改变，两个状态不再属于同一个线性
边值问题，不能直接使用倒易关系。

第一种典型用法是交换“点电荷”和“测量电势的位置”。在导体球外放点电荷时，球面感应电荷分布很复杂；
但把球改成带单位电荷、去掉外部点电荷后，球面电荷均匀，任意外点电势都容易写出。倒易定理把第二种状态
的已知电势转换成第一种状态中整个球的电势，从而避开求感应电荷密度的繁重计算。

第二种用法是交换两个导体。状态 I 中让导体 1 带电 $Q$、导体 2 中性；状态 II 中让导体 2 带相同电荷
$Q$、导体 1 中性。代入倒易式可知，第一状态中导体 2 的电势等于第二状态中导体 1 的电势。这个结论与
两导体的形状是否相同、距离多远都无关，只要求几何在交换状态时没有改变。

第三种用法涉及接地感应电荷。状态 I 中把导体 1 提升到单位电势而让导体 2 接地，状态 II 中交换两导体。
倒易定理说明，两次实验中接地导体上出现的总感应电荷相同。用矩阵语言表述，就是电容系数的互易性
$C_{12}=C_{21}$。这一性质使复杂导体系统所需的独立系数数量大幅减少，也是后来网络互易定理的静电学原型。

倒易定理还可以和能量观点联系。导体组总静电能可写成
\[
 W=\frac12\sum_iQ_i\varphi_i
   =\frac12\sum_{ij}C_{ij}\varphi_i\varphi_j.
\]
交换两个电势变量不应改变同一个能量函数的混合二阶变化，因此系数矩阵必须对称。这给出了倒易关系的另一
种直观来源：互易性并非偶然的代数巧合，而是线性系统存在单值能量函数的结果。

\subsection*{从实验现象到介质极化}

插入电介质的实验可以有两种做法。若电容器先充电后与电源断开，自由电荷 $Q$ 不变，插入介质后电压表读数
降低；由 $C=Q/U$ 可知电容增大。若电容器始终与恒压电源相连，电压不变，插入介质时电源继续向极板输送
电荷；由 $Q=CU$ 可知储荷量增加。两种实验观测的是同一事实，只是外部保持不变的量不同。

同一种介质填满不同形状的电容器时，电容近似按同一比例增加，这一点非常关键。它说明介质作用不能只用
某种特定极板几何来解释，而应由材料自身的响应参数描述。相对介电常数正是这样的宏观参数。表 32.1 中
不同物质的数值相差很大，水的相对介电常数远大于干燥空气，表明分子结构和极化机制会显著影响宏观响应。
击穿场强则描述另一件事：介质在多强的场下还能维持绝缘。介电常数大并不意味着击穿场强必然大，设计真实
电容器时必须同时考虑储能能力和绝缘可靠性。

在保持自由电荷不变的实验中，插入介质后宏观场减小。微观上，外场使原子或分子的正负电荷中心分离，
形成与外场同向的电偶极矩；大量偶极子在介质两端留下束缚电荷。束缚电荷产生的场在介质内部通常与原外场
反向，于是合场减弱。这里“偶极矩与外场同向”和“偶极子附加场与外场反向”并不矛盾：前者描述单个偶极
矩的方向，后者描述极化介质端面束缚电荷在材料内部产生的退极化场。

均匀极化时介质内部的束缚体电荷为零，但表面束缚电荷一般不为零；非均匀极化时 $\rho_b=-\nabla\cdot\vect P$
也会贡献电场。因此，不能把极化简单理解为“只有表面有电荷”。具体是表面项还是体积项占主导，取决于
$\vect P$ 的空间分布和边界形状。

引入 $\vect D$ 的主要好处，是在已知自由电荷而束缚电荷难以逐个计算时直接使用 Gauss 定律。例如平行板
间完全填充均匀线性介质，极板自由电荷面密度为 $\sigma_f$，则 $D=\sigma_f$，进一步由
$D=\varepsilon E$ 得 $E=\sigma_f/\varepsilon$。相对于真空，场强减小 $\varepsilon_r$ 倍，电势差也减小
相同倍数，电容因而增大 $\varepsilon_r$ 倍。这个推导把实验规律、极化图像和宏观场方程连接起来。

不过 $\vect D$ 不能替代 $\vect E$ 描述带电粒子的受力。试探电荷受到的仍是 $q\vect E$，不是
$q\vect D$。在介质界面上，$\vect D$ 的法向跳变由自由面电荷决定，$\vect E$ 的切向分量由静电场保守性
保证连续；只有把两个条件结合起来，才能正确判断电场线如何折转。

\subsection*{线性模型的适用范围}

$\vect P=\varepsilon_0\chi_e\vect E$ 是近似的材料关系。场强过大时，极化可能不再与场成正比；极性分子
的取向还会受到热运动影响，因此介电常数可能随温度变化。交变电场频率很高时，较慢的极化机制来不及响应，
介电常数会随频率变化并出现能量损耗。各向异性晶体中，$\vect P$ 也未必与 $\vect E$ 平行，此时介电常数
应当用张量而不是一个标量表示。

原子极化率 $\alpha$ 是微观量，介电常数是宏观量。若简单地认为每个原子所见局域场就是宏观平均场，便有
$\vect P=N\alpha\vect E$；但邻近偶极子也会改变原子所在位置的场，所以精确关系需要局域场修正。经典的
Clausius--Mossotti 关系写成
\begin{equation}
  \boxed{\frac{\varepsilon_r-1}{\varepsilon_r+2}
  =\frac{N\alpha}{3\varepsilon_0}},
\end{equation}
它展示了如何由宏观介电常数估计微观极化率。该式适用于由近似独立、各向同性分子组成的稀疏介质；对于
强关联固体、极性液体或接近相变的材料，不能机械套用。


\subsection*{重点知识}

\begin{enumerate}[label=\textbf{32.\arabic*},leftmargin=4.2em]
  \item 双极电容器的电容为两极电荷量的绝对值与电势差之比，$C=Q/U$。
  \item 孤立导体的电容为导体电荷量与相对无穷远电势之比，可视为它与无穷远之间的电容。
  \item 电容器储能 $W=Q^2/(2C)=QU/2=CU^2/2$，能量以电场形式存在。
  \item 绝缘体称为电介质。电介质会极化，但一般不会像导体那样形成静电屏蔽和内部等势状态。
  \item 电位移满足 $\vect D=\varepsilon_0\vect E+\vect P$；其 Gauss 定律只以自由电荷为源。
\end{enumerate}
